\subsection{Exact peripheral lifting across abelian chief factors}

All conjugations in this subsection act on the left, and
$[a,b]=aba^{-1}b^{-1}$.  The three peripheral elements of the actual
free profinite group are
\[
 \widehat F=\langle x_1,x_2\rangle,\qquad x_3=(x_1x_2)^{-1},
 \qquad S_i=\overline{\langle x_i\rangle}\simeq\widehat\Z.
 \tag{D.23.1}
\]
In particular the group in the following calculation is the full
profinite group, with these specified three embeddings of
$\widehat\Z$.  For $x_3$ this follows from the explicit basis
replacement $(x_1,x_2)\mapsto(x_1,(x_1x_2)^{-1})$, whose inverse
recovers $x_2=x_1^{-1}x_3^{-1}$.

\begin{statement}{D.23}{Proved}{The complete abelian peripheral obstruction and chief-step count}
Let
\[
 1\longrightarrow A\longrightarrow E\xrightarrow{\pi}G
 \longrightarrow1
 \tag{D.23.2}
\]
be an extension of finite groups with $A$ elementary abelian of
characteristic $p$.  Choose generators $X_1,X_2$ of $E$, put
$X_3=(X_1X_2)^{-1}$ and $g_i=\pi(X_i)$, and denote the resulting
epimorphisms by $q_E:\widehat F\twoheadrightarrow E$ and
$q_G=\pi q_E$.  Let $\alpha\in\operatorname{Aut}(G)$ and
$\lambda\in\widehat\Z^\times$ satisfy
$\alpha(g_i)\sim_G g_i^\lambda$ for all three indices.

Write $A_G=A/[E,A]$, use additive notation in $A$, and set
$C_i=C_G(g_i)$.  There are well-defined homomorphisms
\[
 \kappa_i:C_i\longrightarrow A_G,
 \qquad c\longmapsto[\widetilde c,X_i]\bmod[E,A],
 \qquad H=\sum_{i=1}^3\operatorname{im}\kappa_i.
 \tag{D.23.3}
\]
Choose $h_i\in G$ with
$\alpha(g_i)=h_i g_i^\lambda h_i^{-1}$, choose lifts
$H_i\in E$, and put
\[
 U_i=H_iX_i^\lambda H_i^{-1},\qquad r=U_1U_2U_3\in A.
 \tag{D.23.4}
\]
The class $o(\alpha,\lambda)=[r]\in A_G/H$ is independent of
these choices.  It vanishes exactly when there is a continuous
homomorphism $f:\widehat F\to E$ such that
\[
 \pi f=\alpha q_G,
 \qquad f(x_i)\sim_E X_i^\lambda\quad(1\leq i\leq3).
 \tag{D.23.5}
\]
The number of these literal homomorphisms is zero for nonzero
$o(\alpha,\lambda)$; for zero obstruction it is
\[
 C(E,G;X_1,X_2)=
 \frac{|A|^2|A_G|\prod_{i=1}^3|C_G(g_i)|}
 {|H|\prod_{i=1}^3|C_E(X_i)|}.
 \tag{D.23.6}
\]

When $A$ is an irreducible $\mathbf F_p[G]$-module, every nonempty
set in \textup{(D.23.5)} contains an epimorphism.  More precisely, let
$b(E,G;X_1,X_2)$ be the number of complements $C$ to $A$ in $E$
whose unique section $s_C:G\to C$ satisfies
$s_C(g_i)\sim_E X_i$ for every $i$.  For zero obstruction the number
of epimorphisms in \textup{(D.23.5)} is exactly $C-b>0$.
Thus every noncentral elementary-abelian chief step admits a
generating peripheral lift.  The possible abelian chief obstruction
is confined to a central kernel of order $p$.

The fixed-$h_i$ obstruction $[r]\in A_G$ is the relative extension
class under the explicit identification
\[
 H^2(\widehat F,\{S_1,S_2,S_3\};A_{\alpha q_G})
 \simeq A_G.
 \tag{D.23.7}
\]
Moreover $[r]$ comes from $H_2(G,\mathbf F_p)$ by the universal
central $p$-cover comparison in \textup{(D.23.12)} below.  In
particular $p\nmid|G|$ forces the obstruction to vanish.
\end{statement}

\begin{proof}
Let $\sigma_i$ be conjugation by $U_i$ on $A$.  Its value depends
only on $\alpha(g_i)$, and
$\sigma_1\sigma_2\sigma_3=1$.  Replacing $H_i$ by $a_iH_i$, for
$a_i\in A$, changes the product in \textup{(D.23.4)} to
\[
 r+L(a_1,a_2,a_3),\qquad
 L(a_1,a_2,a_3)=(1-\sigma_1)a_1
       +\sigma_1(1-\sigma_2)a_2
       +\sigma_1\sigma_2(1-\sigma_3)a_3.
 \tag{D.23.8}
\]
This is a multiplication identity in the extension, since all three
correction factors lie in the abelian normal subgroup $A$.
The sum of the first two image spaces is
\[
 (1-\sigma_1)A+(1-\sigma_1\sigma_2)A=[E,A].
\]
For the equality, use $1-ab=(1-a)b+(1-b)$ and the fact that
$\sigma_1,\sigma_1\sigma_2$ generate the action of $G$; inverses
give the same image spaces.  The third image space is already in
this sum.  Consequently $\operatorname{im}L=[E,A]$, and
$|\ker L|=|A|^2|A_G|$.  This proves the assertion with the three
base conjugators fixed.

To vary those conjugators, work in $\overline E=E/[E,A]$.
Its kernel $A_G$ is central.  If $c\in C_i$, the commutator
$[\widetilde c,X_i]$ lies in $A$; changing the lift of $c$ changes
its image in $A_G$ by zero.  The central commutator identity
$[ab,X_i]=[a,X_i][b,X_i]$ proves that $\kappa_i$ is a homomorphism.
All conjugators in $G$ from $g_i^\lambda$ to $\alpha(g_i)$ have
the form $h_ic_i$, with $c_i\in C_i$, since raising to the unit
$\lambda$ preserves the generated cyclic subgroup.  Their effect
on the product class is
\[
 [r]\longmapsto[r]+\lambda\sum_i\kappa_i(c_i).
 \tag{D.23.9}
\]
Here conjugation by a lift of $h_i$ acts trivially on $A_G$ and
$[\widetilde c,X_i^\lambda]=[\widetilde c,X_i]^\lambda$ in
$\overline E$.  Multiplication by $\lambda$ is an automorphism of
$A_G$, and preserves $H$.  This proves choice independence and the
criterion in $A_G/H$.

For the count, a zero obstruction gives exactly
$\prod_i|C_i|/|H|$ triples of base conjugators satisfying the
coinvariant equation.  Each such triple has exactly
$|\ker L|=|A|^2|A_G|$ triples of lifts in $E$ whose conjugated
peripheral product is $1$.  Each resulting triple of peripheral
elements is counted $\prod_i|C_E(X_i^\lambda)|$
times.  These centralizers equal $C_E(X_i)$ because $\lambda$ is a
unit.  A triple with product $1$ gives exactly one continuous
homomorphism out of $\widehat F$, by its free profinite universal
property.  This proves \textup{(D.23.6)} without quotienting by a
centralizer that was not counted.

Now assume $A$ irreducible.  If $L=f(\widehat F)$ for a solution,
then $L\to G$ is onto and $L\cap A$ is $G$-invariant: conjugation
by any element of $G$ can be computed using a lift in $L$.
Hence $L\cap A$ is either $A$ or $0$.  The first case gives $L=E$;
the second makes $L$ a complement.  For a given complement $C$,
there is exactly one homomorphism over $\alpha q_G$ with image
in $C$, namely $s_C\alpha q_G$.  It satisfies the peripheral
conditions precisely when $s_C(g_i)\sim_E X_i$ for all $i$.
Indeed
\[
 s_C(\alpha(g_i))=
 s_C(h_i)s_C(g_i)^\lambda s_C(h_i)^{-1},
\]
and taking the power $\lambda^{-1}$ in the finite group $E$
shows that $s_C(g_i)^\lambda\sim_E X_i^\lambda$ is equivalent
to $s_C(g_i)\sim_E X_i$.  Thus the nongenerating solutions are
in a literal bijection with the same $b$ admissible complements,
independently of $\alpha$ and $\lambda$.  The identity case
$(\alpha,\lambda)=(1,1)$ contains the given epimorphism $q_E$,
so its count $C$ is strictly greater than $b$.  Every other
nonempty case has this same count $C$, proving the claimed
generating-lift theorem.  Finally an irreducible module with
$A_G\ne0$ is the one-dimensional trivial module: the quotient
map to a nonzero trivial quotient has zero kernel by irreducibility.
This proves the noncentral chief assertion.

For clarity, the count in \textup{(D.23.6)} is for homomorphisms,
with the specified base map retained.  If one further takes
$A$-conjugacy classes of the generating solutions, every orbit
has size $|A|/|A^G|$.  If one permits all conjugations in $E$
that retain the same base map $\alpha q_G$, the acting group is
$\pi^{-1}Z(G)$ and every generating stabilizer is $Z(E)$.
The corresponding number of classes is therefore
$(C-b)|Z(E)|/(|A||Z(G)|)$.  These formulas specify the extra
quotient maps rather than silently identifying literal lifts with
their classes.

Here is the relative-cohomology comparison.  The free profinite
universal property makes the cocycle group
$Z^1(\widehat F,A)$ equal to $A^2$: arbitrary values on $x_1,x_2$
define a continuous homomorphism to the finite semidirect product
associated with the action.  Each $S_i\simeq\widehat\Z$ has
$H^1(S_i,A)=A/(\sigma_i-1)A$.  Projectivity of $\widehat F$
gives $H^2(\widehat F,A)=0$, so the relative exact sequence
identifies the group in \textup{(D.23.7)} with the cokernel of
the restriction map on $H^1$.  On cocycle values the latter is
\[
 (u,v)\longmapsto
       (u,v,-\sigma_3(u+\sigma_1v)).
 \tag{D.23.10}
\]
The map from the three local values to this cokernel, followed
by \textup{(D.23.7)}, is
\[
 (b_1,b_2,b_3)\longmapsto
       b_1+\sigma_1b_2+\sigma_1\sigma_2b_3\pmod{[E,A]}.
 \tag{D.23.11}
\]
It is onto.  Eliminating $b_1,b_2$ using
\textup{(D.23.10)} leaves exactly the relations
$(\sigma_1-1)A+\sigma_1(\sigma_2-1)A=[E,A]$, which proves
that its kernel is the asserted restriction image.
Pull back \textup{(D.23.2)} along $\alpha q_G$ and prescribe the
local sections that send $x_i$ to $(U_i,x_i)$.  Choose the global
splitting taking $x_1,x_2$ to $(U_1,x_1),(U_2,x_2)$.
The difference of the third prescribed section from this splitting
has cocycle value $\sigma_3r$, so
\textup{(D.23.11)} sends the relative class to $r$.
This is the actual finite-kernel relative extension of the
full profinite pair.  It agrees with the classification and
lifting comparison of Wilkes, Theorem~5.5 and Proposition~6.3,
with lifting problems defined in Definition~6.1
\cite{APeripheralWilkes2025}.  In particular no passage from a
pro-$p$ pair to the full profinite pair is being assumed.

For the final assertion put $K=\ker q_G$ and
$U=\widehat F/\overline{K^p[\widehat F,K]}$.
Finite-coset Schreier rewriting gives at most $1+|G|$ topological
generators for $K$.  Therefore its elementary abelian quotient
$A_U$ is finite, and it is central in $U$.  The homology five-term sequence
and $H_2(\widehat F,\mathbf F_p)=0$ give the exact comparison
\[
 0\longrightarrow H_2(G,\mathbf F_p)\longrightarrow
 A_U=K/\overline{K^p[\widehat F,K]}
 \longrightarrow H_1(\widehat F,\mathbf F_p)
 \longrightarrow H_1(G,\mathbf F_p)\longrightarrow0.
 \tag{D.23.12}
\]
The original map $q_E$ induces an actual surjection
$U\twoheadrightarrow E/[E,A]$, and its map on kernels is
$A_U\twoheadrightarrow A_G$.  Choose lifts $\widetilde H_i$
of the $H_i$ in $\widehat F$.  The word
\[
 \prod_{i=1}^3\widetilde H_i x_i^\lambda
                    \widetilde H_i^{-1}
 \tag{D.23.13}
\]
lies in $K$ and has exactly zero image in
$\widehat F^{\rm ab}=\widehat\Z^2$, since
$\lambda(x_1+x_2+x_3)=0$ there.  Its image in $A_G$ is $[r]$.
Thus its class in $A_U$ belongs to the $H_2(G,\mathbf F_p)$
term of \textup{(D.23.12)}.  If $p\nmid|G|$, averaging contracts
the positive-degree finite-group homology complex over
$\mathbf F_p$, so that term is zero.  For general $G$ the Schur
term and the quotient by the three maps \textup{(D.23.3)} both
remain in the calculation.
\end{proof}

The independent finite program \texttt{a\_peripheral\_count.py}
enumerates actual peripheral triples and their generated subgroups,
then compares them with \textup{(D.23.6)} and the complement
bijection.  It checks sixteen cases in six extensions.  In the
$V_4\triangleleft A_4$ example there are sixteen weak solutions,
twelve generating solutions, and four admissible complements.  In
the central $C_2\triangleleft D_8$ example the centralizer-image
subgroup $H$ has order two; including it gives the four actual
solutions.  These examples test both terms that an unrestricted
generator-lifting argument would omit.  The script, exact output,
and JSON tables are retained in \texttt{code/}.

\subsection{A nonzero spin obstruction and its exact pro-$\ell$ fibre}

\begin{statement}{D.24}{Proved}{A peripheral-preserving finite automorphism with no peripheral lift}
Let $S=\operatorname{PSL}_2(\mathbf F_{11})$ and
\[
 E=\operatorname{SL}_2(\mathbf F_{11})\times S,
 \qquad G=S\times S,
 \qquad\pi=(\text{central quotient},1).
 \tag{D.24.1}
\]
There is an explicit epimorphism $q_E:\widehat F\twoheadrightarrow E$
for which the factor-swap automorphism
$\alpha:G\to G$, $(a,b)\mapsto(b,a)$, preserves each of the
three labelled peripheral conjugacy classes with multiplier one,
but no homomorphism \textup{(D.23.5)} exists.  Its obstruction is
the nonidentity element of $\ker\pi=\{(I,1),(-I,1)\}$.

The maximal pro-$\ell$ quotients of both $E$ and $G$ are trivial
for every prime $\ell$.  More precisely, for any two-generated
finite $\ell$-group $L$ with a specified generating triple
$l_1l_2l_3=1$ and any multiplier-one peripheral automorphism
$\beta\in\operatorname{Aut}(L)$, the corresponding product
quotient has maximal pro-$\ell$ quotient exactly $L$.
The automorphism $\alpha\times\beta$ of $G\times L$ has no
peripheral lift to $E\times L$, whereas $1_G\times\beta$ has
an explicit generating peripheral lift.  Thus the same prescribed
finite pro-$\ell$ component has both an obstructed and an admitted
choice of perfect component in this example.
\end{statement}

\begin{proof}
All matrix entries in the following display are in $\mathbf F_{11}$.
Take the two triples of canonical odd-order lifts
\[
\begin{array}{c|ccc}
 i&1&2&3\\\hline
 A_i&
 \begin{pmatrix}0&10\\1&10\end{pmatrix}&
 \begin{pmatrix}0&8\\4&2\end{pmatrix}&
 \begin{pmatrix}6&2\\4&7\end{pmatrix}\\[3pt]
 B_i&
 \begin{pmatrix}0&10\\1&10\end{pmatrix}&
 \begin{pmatrix}1&0\\1&1\end{pmatrix}&
 \begin{pmatrix}1&10\\0&1\end{pmatrix}.
\end{array}
 \tag{D.24.2}
\]
Write $a_i,b_i$ for their respective images in $S$.
Direct multiplication and powering give
\[
 A_1A_2A_3=I,\qquad B_1B_2B_3=-I,
 \qquad |A_i|=|B_i|=(3,11,11)_i.
 \tag{D.24.3}
\]
The explicit matrices
\[
 H_1=\begin{pmatrix}0&1\\10&1\end{pmatrix},\quad
 H_2=\begin{pmatrix}2&6\\0&6\end{pmatrix},\quad
 H_3=\begin{pmatrix}0&5\\2&3\end{pmatrix}
 \tag{D.24.4}
\]
have determinant one and satisfy $H_iA_iH_i^{-1}=B_i$.
In particular $a_i$ and $b_i$ are conjugate in $S$ for each $i$.
Set
\[
 X_i=(A_i,b_i)\in E,\qquad g_i=(a_i,b_i)\in G.
 \tag{D.24.5}
\]
Both triples have product one.  They generate their displayed
groups, as certified by the complete matrix enumeration below.
The elements $(\overline H_i,\overline H_i^{-1})\in G$
conjugate $g_i$ to $\alpha(g_i)$, verifying the exact base
peripheral condition.

An element of $E$ above $\alpha(g_i)=(b_i,a_i)$ has first
coordinate either $B_i$ or $-B_i$.  Exactly one of these has odd
order, namely $B_i$: if $B_i$ has odd order $n$, then
$(-B_i)^n=-I$ and $-B_i$ has order $2n$.  Every conjugate of
$X_i$ has the same odd order as $X_i$, so its only possible
value above $\alpha(g_i)$ is
\[
 U_i=(B_i,a_i).
 \tag{D.24.6}
\]
It is indeed conjugate to $X_i$ in $E$, using
$(H_i,\overline H_i^{-1})$.  But its product is
$U_1U_2U_3=(-I,1)\ne1$.  Thus no choice of conjugators, and no
nongenerating homomorphism, solves this peripheral lifting problem.
Equivalently, each base-centralizer homomorphism
\textup{(D.23.3)} is zero.  To see this directly, in this central
extension the commutator $[\widetilde c,X_i]$ has order dividing
two and also dividing the odd order of $X_i$, so it is one.
The computed nonzero relative class consequently survives the
full base-centralizer quotient.

Here is the finite certificate for all generation assertions.
The program \texttt{a\_peripheral\_spin.py} constructs all
determinant-one $2\times2$ matrices over $\mathbf F_{11}$ and
identifies a matrix with its negative for $S$.  It verifies their
orders $1320$ and $660$, the three canonical odd-order lifts,
all products in \textup{(D.24.3)}, and all three conjugations in
\textup{(D.24.4)}.  Starting from the identity, it repeatedly
right-multiplies the generated set by the two displayed generators
until the set stops growing.  This computes the generated subgroup
exactly: the set contains every positive word, is closed under
multiplication by the generators, and positive powers include
their inverses in these finite groups.  For $(g_1,g_2)$ it has
$435600=660^2$ elements, and for $(X_1,X_2)$ it has
$871200=1320\cdot660$ elements.  These are the complete group
orders, proving surjectivity.  The matrices, canonical lifts,
products, conjugators, and both cardinalities are saved in
\texttt{a\_peripheral\_spin\_p11.json}; the executed output is
\texttt{a\_peripheral\_spin\_p11.out}.  The separately retained
$p=5,7$ searches found no matching pair in their stated exhaustive
odd-order search domains; the proof here uses the displayed
$p=11$ certificate.  The program explicitly restricts its finite
order bound to these three tested primes.

For the pro-$\ell$ statement, first prove perfectness directly.
In $\operatorname{SL}_2(\mathbf F_{11})$ put
\[
 D=\operatorname{diag}(2,6),\quad
 u(t)=\begin{pmatrix}1&t\\0&1\end{pmatrix},\quad
 v(t)=\begin{pmatrix}1&0\\t&1\end{pmatrix}.
\]
One has $[D,u(t)]=u(3t)$ and $[D,v(t)]=v(2t)$.
Both scalar factors are invertible in $\mathbf F_{11}$.
Elementary row reduction expresses every determinant-one matrix
as a product of the $u(t),v(t)$, so the whole group is its
commutator subgroup.  Its quotient $S$ is perfect, and hence
so are $E$ and $G$.  A finite perfect group has trivial image
in every finite $\ell$-group: such an image would be a perfect
finite $\ell$-group, whereas induction on the central series
shows that a nontrivial finite $\ell$-group has nontrivial
abelianization.  Passing to the finite quotients of a pro-$\ell$
group proves the assertion about maximal pro-$\ell$ quotients.

Finally use the source generators $(X_i,l_i)$ and $(g_i,l_i)$
in $E\times L$ and $G\times L$.  They generate the full products.
Indeed their generated subgroups project onto both factors; a
subdirect subgroup $B\leq P\times L$, with coordinate kernels
$P_0=\{p:(p,1)\in B\}$ and $L_0=\{l:(1,l)\in B\}$, gives
the isomorphism $P/P_0\to L/L_0$, $pP_0\mapsto lL_0$ for
$(p,l)\in B$.  For perfect
$P$ and an $\ell$-group $L$, this common quotient is perfect
and solvable, hence trivial, which makes the subgroup the full
product.  This also proves that the actual coordinate projection
is the maximal pro-$\ell$ quotient map of either product.
The shadow $\alpha\times\beta$ satisfies the peripheral
conditions by the separate conjugators in the two factors.
Any lift would project to the already excluded lift of $\alpha$
in \textup{(D.24.6)}.  In contrast, the literal triple
\[
 (X_i,\beta(l_i))\quad(1\leq i\leq3)
 \tag{D.24.7}
\]
has product one, has the prescribed three conjugacy classes,
and generates $E\times L$ by the same subdirect-product
argument.  It is the promised lift of $1_G\times\beta$.
\end{proof}

Statements D.23--D.24 supply an exact chief-step obstruction and
a nonzero full-profinite peripheral example.  The displayed
maximal pro-$\ell$ quotient maps and the two product lifts specify
their effect on the pro-$\ell$ comparison.  They do not choose
between the dimensions in D.1 or D.2: the choice of compatible
finite components, including coupled mixed quotients, remains
part of those two problems.
